% BHU PYQ -- Hindi: पद्य एवं संप्रेषण - कविता (2024-25 NEP AEC)
\documentclass[11pt,a4paper]{article}
\usepackage[a4paper,top=2.5cm,bottom=2.5cm,left=2.8cm,right=2.8cm,headheight=14pt]{geometry}
\usepackage{amsmath,amssymb}
\usepackage{fontspec}
\setmainfont{Kohinoor Devanagari}
\usepackage{microtype,fancyhdr,enumitem,hyperref,lastpage,parskip}

\hypersetup{colorlinks=true,urlcolor=black,linkcolor=black,pdftitle={HINAE11: पद्य एवं संप्रेषण - कविता -- BHU NEP 2024-25}}

\pagestyle{fancy}\fancyhf{}
\renewcommand{\headrulewidth}{0.4pt}\renewcommand{\footrulewidth}{0.4pt}
\fancyhead[L]{\small \textit{Banaras Hindu University}}
\fancyhead[R]{\small \textit{HINAE11: पद्य एवं संप्रेषण - कविता (2024-25)}}
\fancyfoot[C]{\small Page \thepage\ of \pageref{LastPage}}

\newcommand{\pts}[1]{\hfill \textit{[#1]}}

\begin{document}

\begin{center}
  {\large\bfseries BANARAS HINDU UNIVERSITY}\\[2pt]
  {\normalsize B.A. (Hons.) Semester-I Examination, 2024-25 (NEP)}\\[6pt]
  \rule{0.8\linewidth}{0.5pt}\\[6pt]
  {\large\bfseries HINDI}\\[4pt]
  {\large\bfseries Paper Code: HINAE11 : पद्य एवं संप्रेषण - कविता}\\[4pt]
  {\normalsize Time: 2:30 Hours \hfill Full Marks: 70}\\[2pt]
  {\small REF NO. 2024-25/D-I/091}
\end{center}
\rule{\linewidth}{0.4pt}
\smallskip


\noindent \textit{सभी प्रश्नों के उत्तर दीजिए।}

\bigskip

\section*{खण्ड -- अ / अति लघुउत्तरीय प्रश्न}
\noindent \textit{नोट: किन्हीं \textbf{पाँच} प्रश्नों के उत्तर दीजिए।}\pts{5 $\times$ 4 = 20}

\begin{enumerate}[label=\textbf{(\ क\ )},leftmargin=*,itemsep=6pt]
  \item[\textbf{(क)}] पाठ्य विषय की प्रमुख प्रवृत्तियों का उल्लेख कीजिए।\pts{4}
  \item[\textbf{(ख)}] भाषा और शैलीगत विशेषताओं की व्याख्या कीजिए।\pts{4}
  \item[\textbf{(ग)}] प्रमुख रचनाकारों के योगदान पर टिप्पणी लिखिए।\pts{4}
  \item[\textbf{(घ)}] विषय के सामाजिक-सांस्कृतिक महत्त्व को स्पष्ट कीजिए।\pts{4}
  \item[\textbf{(ङ)}] व्यावहारिक अनुप्रयोग पर प्रकाश डालिए।\pts{4}
\end{enumerate}

\bigskip

\section*{खण्ड -- ब / दीर्घउत्तरीय प्रश्न}
\noindent \textit{नोट: किन्हीं \textbf{तीन} प्रश्नों के उत्तर दीजिए।}\pts{3 $\times$ 10 = 30}

\begin{enumerate}[label=\textbf{\arabic*.},leftmargin=*,itemsep=12pt]
  \item विषयवस्तु और प्रासंगिकता का सोदाहरण विश्लेषण कीजिए।\pts{10}
  \item प्रमुख सिद्धांतों एवं प्रवृत्तियों की समीक्षा कीजिए।\pts{10}
  \item आधुनिक संदर्भों में इसके महत्त्व का निरूपण कीजिए।\pts{10}
\end{enumerate}

\bigskip

\section*{खण्ड -- स / व्याख्या एवं टिप्पणियाँ}
\noindent \textit{नोट: किन्हीं \textbf{दो} अंशों की सप्रसंग व्याख्या / टिप्पणी कीजिए।}\pts{2 $\times$ 10 = 20}

\begin{enumerate}[label=\textbf{(\ क\ )},leftmargin=*,itemsep=10pt]
  \item[\textbf{(क)}] प्रमुख काव्यांश / पाठ्यांश की सप्रसंग व्याख्या कीजिए।\pts{10}
  \item[\textbf{(ख)}] प्रासंगिक गद्यांश / काव्यांश का भावार्थ स्पष्ट कीजिए।\pts{10}
\end{enumerate}


\end{document}
